\chapter{ Positional and Keyword Arguments}


\section{Introduction}
In the previous sessions, you wrote several \textbf{functions} and called them from a \texttt{main()} function.
That already gave you a powerful skill: you can break a larger program into smaller, reusable pieces.

In Session 5 we keep that same idea, but we make your functions \textbf{more flexible} by controlling them through
\textbf{arguments} when you call them.

\noindent
\textbf{Parameters vs.\ arguments (a common beginner confusion).}
\begin{itemize}\setlength{\itemsep}{2pt}
    \item A \textbf{parameter} is the name written in the function definition (the ``slot'' the function expects).
    \item An \textbf{argument} is the actual value you provide when calling the function (what gets put into the slot).
\end{itemize}

\noindent
Example:
\begin{minted}[fontsize=\small, linenos]{python}
def greet(name, punctuation="!"):   # name and punctuation are parameters
    print("Hello", name + punctuation)

greet("Ada")                        # "Ada" is an argument
greet(name="Ada", punctuation="?")  # both are arguments (keyword style)
\end{minted}

\noindent
In this session's tutorial, we focus on \textbf{two} ways of providing arguments:
\begin{itemize}\setlength{\itemsep}{2pt}
    \item \textbf{Positional arguments:} values are matched to parameters by \textit{order}.
    \item \textbf{Keyword arguments:} values are matched to parameters by \textit{name}.
\end{itemize}

Understanding these techniques matters because they directly affect:
\begin{itemize}\setlength{\itemsep}{2pt}
    \item how easy your code is to read (especially for someone else),
    \item how easy it is to change your code later (refactoring without breaking calls),
    \item how likely you are to introduce subtle bugs (wrong value in the wrong slot).
\end{itemize}



\section{What You Will Learn:}

Upon completing this tutorial, you'll be able to:
\begin{itemize}\setlength{\itemsep}{0pt}
    \item \textbf{Positional arguments:} How Python matches arguments to parameters by position, and why the order matters.
    \item \textbf{Keyword arguments:} How naming parameters in a call improves readability and lets you override only what you need.
    \item \textbf{Default values:} How ``optional'' parameters work, how \texttt{None} is commonly used as a ``not provided'' marker, and how defaults interact with positional and keyword calls.
    \item \textbf{Beginner-safe debugging:} How to recognize and fix common argument errors such as ``missing required argument'' or ``multiple values for argument''.
\end{itemize}


\noindent
\textbf{Important framing:} Session 5 does not replace the session 4 classifier.
Instead, it extends the same function set so you can \textit{control the behavior} through function arguments
instead of rewriting the classifier logic each time you want a small change.
This is a key refactoring mindset: \textbf{keep the core rule stable, make the interface flexible}.

\section{Learning Goal}
By the end of Session 5, you will:
\begin{enumerate}
    \item keep session 4 function names and core classifier behavior unchanged,
    \item add positional parameters so a caller can override key values \textit{in order},
    \item add keyword parameters with sensible defaults so a caller can override values \textit{by name},
    \item learn how to import a helper function from another file (your first step toward multi-file projects),
    \item call your pipeline like a small API from \texttt{main()} to demonstrate different configurations clearly.
\end{enumerate}
\noindent
If you achieve these goals, you will have a program that can run the same pipeline in multiple modes without duplicating code.

\subsection{Hands-on Files for This Session}
Use:
\begin{itemize}
    \item \texttt{session5/session5\_iris\_template.py} (your main working file for this session)
    \item \texttt{session5/iris\_data.csv} (the dataset you will load and classify)
    \item \texttt{session5/helper/iris\_utils.py} (provided helper: a simple ``fit'' function for learning a threshold)
\end{itemize}
\noindent
\textbf{What you are starting with:} a mostly-complete session 5 template where the core classifier functions already exist.
\textbf{What you will modify:} function signatures and call sites so the pipeline can be configured using arguments.

\section{Importing Packages and Functions}
In Python, as your project grows, you will want to split your code into multiple files. You can import functions from one file into another using the \texttt{import} statement.

For this session, we provide a helper module \texttt{session5/helper/iris\_utils.py} which contains a function to calculate a threshold from the dataset.
You will \textbf{not} rewrite this helper; you will \textbf{use} it.
Import it at the top of your script like this:

\begin{minted}[fontsize=\small, linenos]{python}
import csv
from helper.iris_utils import fit_threshold_from_setosa
\end{minted}

\noindent
\textbf{Why this import is important for beginners:}
\begin{itemize}
    \item In Session 4, all logic stayed in one file.
    \item In Session 5, we start separating logic across files so code is easier to maintain.
    \item \texttt{helper.iris\_utils} means: ``go to the \texttt{helper} folder and use the function inside \texttt{iris\_utils.py}''.
\end{itemize}
\noindent
\textbf{Common misunderstanding:} beginners often think the import path is the same as the file path.
In Python, the import path depends on \textit{where you run the script from}.
For this tutorial, run from the project root:
\begin{minted}[fontsize=\small, linenos]{bash}
python3 session5/session5_iris_template.py
\end{minted}
This makes \texttt{helper} importable because the script is located inside \texttt{session5/}.

\subsection{Session 4 vs Session 5 (Function-by-Function Changes)}
Session 5 is a refactor focused on \textbf{interfaces}, not new logic.
You should be able to point to each change and say:
\begin{quote}
``I did this so callers can control the pipeline through arguments instead of editing code.''
\end{quote}

\noindent
Use this checklist while coding to stay on track:
\begin{itemize}
    \item \texttt{load\_iris\_data(filepath)} in Session 4 $\rightarrow$ \texttt{load\_iris\_data(filepath=None)} in Session 5.
    \item \texttt{initialize\_predictions(dataset, settings)} in Session 4 $\rightarrow$ \texttt{initialize\_predictions(dataset, settings, print\_each=True)} in Session 5.
    \item \texttt{setup\_application(filepath)} in Session 4 $\rightarrow$ \texttt{setup\_application(filepath=None, threshold=None, print\_each=True, use\_fit=False)} in Session 5.
    \item \texttt{main()} in Session 4 ran one flow $\rightarrow$ Session 5 runs default, positional, keyword, and fit-based calls.
    \item \textbf{Unchanged core classifier rule:} threshold comparison logic remains exactly the same.
\end{itemize}
\noindent
The key pattern is: \textbf{new parameters} (inputs) and \textbf{defaults} (fallback behavior), while keeping the same prediction rule.

\subsection{Canonical Function Set (Session 5)}
Complete these functions in \texttt{session5\_iris\_template.py}. The table below shows:
\begin{itemize}
    \item the \textbf{function section/name},
    \item the \textbf{code block} to place inside the function,
    \item the \textbf{notes} explaining what changed in Session 5.
\end{itemize}

\input{session5/tab5.tex}

Classifier rule remains:
\begin{minted}[fontsize=\small, linenos]{python}
if sample[feature_name] < threshold:
    y_pred = positive_label
else:
    y_pred = negative_label
\end{minted}
\noindent
\textbf{Keep reminding yourself:} in this session, you are improving \textit{how the functions are called}, not changing \textit{how the classifier decides}.

\subsection{Introduce \texttt{fit()}}
In real machine learning, \texttt{fit()} means ``learn parameters from data'' (for example, learn weights for a model).
That topic is much bigger than this module, but we can introduce the \textbf{idea} in a simple, beginner-friendly way:
we will learn \textbf{one} number (a threshold) from the dataset.

\noindent
\textbf{Why do this?} It teaches a useful design lesson:
\begin{itemize}\setlength{\itemsep}{2pt}
    \item \textbf{Manual configuration:} sometimes you choose parameters yourself (e.g., ``threshold = 2.0'').
    \item \textbf{Data-driven configuration:} sometimes you compute parameters from data (a very small version of ``learning'').
\end{itemize}

At the beginning of this session, we provide the fit helper function as-is in \texttt{session5/helper/iris\_utils.py} (you can reuse it directly):
\begin{minted}[fontsize=\small, linenos]{python}
def fit_threshold_from_setosa(dataset, fallback_threshold=2.0):
    total = 0.0
    count = 0

    for sample in dataset:
        if sample["species"] == "setosa":
            total += sample["petal_length"]
            count += 1

    return total / count if count > 0 else fallback_threshold
\end{minted}

You will then call this through \texttt{setup\_application(..., use\_fit=True)} so the threshold can be learned automatically.
\noindent
\textbf{Important:} ``fit'' here does not change the classifier rule. It only chooses the threshold value that the rule will use.


\subsection{Positional Arguments}
\subsection{Task 1: Positional Arguments (First)}
This task teaches you the most basic argument style: \textbf{positional arguments}.
When you call a function positionally, Python matches each argument to a parameter purely by \textbf{left-to-right order}.

\noindent
\textbf{What you are starting with:} functions that work only with built-in defaults (or require you to edit code to change behavior).

\noindent
\textbf{What you will modify:} function definitions (signatures) so they accept parameters like \texttt{filepath}, \texttt{threshold}, and \texttt{print\_each}.

\noindent
\textbf{Why this improves your code:} you can run experiments (different thresholds, different verbosity) by changing the \textit{call}, not the \textit{function body}.

\noindent
\textbf{What should change afterwards:} calling \texttt{setup\_application(...)} with different positional values should lead to different threshold/printing behavior, while the classifier rule stays the same.

Example call with positional arguments:
\begin{minted}[fontsize=\small, linenos]{python}
settings, dataset, result = setup_application(
    "session5/iris_data.csv",
    1.8,
    False
)
\end{minted}

The positional order must match your function definition.
\noindent
\textbf{Common beginner pitfall:} positional arguments are short, but easy to mix up.
For example, accidentally swapping \texttt{threshold} and \texttt{print\_each} may not crash your program, but it will produce confusing behavior.
This is one reason keyword arguments are often preferred for readability (Task 2).

\subsection{Detailed Step-by-Step Coding Manual (Beginner Version)}
Follow these steps \textbf{in order} inside \texttt{session5/session5\_iris\_template.py}.

\begin{enumerate}
    \item \textbf{Start from Session 4 structure.}
    Keep all Session 4 function names and keep the same classifier rule logic.
    Your goal is not to create new classifier logic; it is to make the existing logic configurable through parameters.

    \item \textbf{Add imports at the very top.}
    \begin{minted}[fontsize=\small, linenos]{python}
    import csv
    from helper.iris_utils import fit_threshold_from_setosa
    \end{minted}
    \noindent
    \textbf{Reasoning:} importing a helper keeps your main file focused on the pipeline, while the helper file contains one reusable ``learning'' function.

    \item \textbf{Update the data-loading signature to support defaults.}
    Change from Session 4 style:
    \begin{minted}[fontsize=\small, linenos]{python}
    def load_iris_data(filepath):
    \end{minted}
    to Session 5 style:
    \begin{minted}[fontsize=\small, linenos]{python}
    def load_iris_data(filepath=None):
        if filepath is None:
            print("[WARNING] No filepath provided. Using default 'session5/iris_data.csv'.")
        filepath = "session5/iris_data.csv"
    \end{minted}
    \noindent
    \textbf{Reasoning:} beginners can run the file without needing to pass arguments immediately, while advanced users can still override the path.

    \item \textbf{Keep CSV parsing logic unchanged.}
    Continue using \texttt{csv.DictReader} and \texttt{float(...)} conversion exactly as in Session 4.
    \noindent
    \textbf{Reasoning:} refactoring is safest when you change one thing at a time. Here, we refactor argument handling, not CSV reading.

    \item \textbf{Keep prediction helper functions unchanged.}
    \texttt{determine\_binary\_label}, \texttt{determine\_true\_binary\_label}, and \texttt{evaluate\_prediction} should keep the same logic from Session 4.
    \noindent
    \textbf{Reasoning:} these functions encode the core classification rule and the evaluation rule, which should remain stable in this session.

    \item \textbf{Add optional output control in initialize\_predictions.}
    Change the function header to:
    \begin{minted}[fontsize=\small, linenos]{python}
    def initialize_predictions(dataset, settings, print_each=True):
    \end{minted}
    Then print per-sample lines \textbf{only if} \texttt{print\_each} is True.
    \noindent
    \textbf{Reasoning:} as your datasets grow, printing every sample becomes noise. A boolean flag lets you keep the option without deleting the print statements.

    \item \textbf{Compute and store accuracy in metrics.}
    Keep the Session 4 formula and write into \texttt{metrics["accuracy"]}.
    \noindent
    \textbf{Expected outcome:} your result dictionary should always include \texttt{"correct"}, \texttt{"wrong"}, \texttt{"total"}, and \texttt{"accuracy"} so later code can print consistent summaries.

    \item \textbf{Upgrade setup\_application into an API-style function.}
    Use this Session 5 header:
    \begin{minted}[fontsize=\small, linenos]{python}
    def setup_application(filepath=None, threshold=None, print_each=True, use_fit=False):
    \end{minted}
    Then apply priority in this order:
    \begin{enumerate}
        \item if \texttt{use\_fit} is True, learn threshold from \texttt{fit\_threshold\_from\_setosa(...)};
        \item else if \texttt{threshold} is provided, use manual threshold;
        \item else keep default threshold from \texttt{get\_default\_settings()}.
    \end{enumerate}
    \noindent
    \textbf{Reasoning:} sometimes you want ``automatic'' configuration (fit), sometimes you want manual configuration (threshold), and sometimes you want to accept defaults.
    The priority rule makes your program predictable. If \texttt{use\_fit=True}, it should be obvious that the learned threshold wins, even if a manual threshold was also provided.

    \item \textbf{Add a summary printer function.}
    Implement \texttt{print\_summary(title, result)} so each run reports Correct/Wrong/Total/Accuracy.
    \noindent
    \textbf{Reasoning:} beginners learn faster when they can see the effect of a change immediately.
    A single summary function also avoids copy-pasting four nearly identical print blocks.

    \item \textbf{Write main() with four required runs.}
    In this exact teaching order:
    \begin{enumerate}
        \item default call,
        \item positional override call,
        \item keyword override call,
        \item fit-based call using \texttt{use\_fit=True}.
    \end{enumerate}
    \noindent
    \textbf{After you run:} you should see four separate summaries. The threshold value used may differ between runs, and the accuracy may change with it.

    \item \textbf{Keep beginner-safe script entry point.}
    \begin{minted}[fontsize=\small, linenos]{python}
    if __name__ == "__main__":
        main()
    \end{minted}

    \item \textbf{Run and verify from project root.}
    Use:
    \begin{minted}[fontsize=\small, linenos]{bash}
    python3 session5/session5_iris_template.py
    \end{minted}
    \noindent
    \textbf{What to check:}
    \begin{itemize}\setlength{\itemsep}{2pt}
        \item The program runs without you editing thresholds inside function bodies.
        \item The positional run and keyword run actually change the threshold or printing behavior.
        \item The fit-based run produces a threshold that comes from the setosa samples (not just the default).
    \end{itemize}
\end{enumerate}

\subsection{Keyword Arguments with Defaults}

\subsection{Task 2: Keyword Arguments with Defaults}
This task teaches you how to call functions in a way that is often \textbf{more readable} than positional arguments.
With keyword arguments, you explicitly name which parameter you are setting.

\noindent
\textbf{Why beginners like keyword arguments:} you do not have to remember the exact order, and your code becomes self-documenting.
Compare these two calls:
\begin{minted}[fontsize=\small, linenos]{python}
setup_application("session5/iris_data.csv", 1.8, False)
setup_application(filepath="session5/iris_data.csv", threshold=1.8, print_each=False)
\end{minted}
Both are valid, but the keyword call makes it harder to accidentally swap values.

\noindent
Make the same parameters optional with defaults (for example \texttt{filepath=None}, \texttt{threshold=None}, \texttt{print\_each=True}).

Inside your function, when a parameter is \texttt{None}, fallback to default values.
\noindent
\textbf{Reasoning:} \texttt{None} is a common Python convention meaning ``no value was provided''.
It lets you distinguish between ``user chose a value'' and ``user did not set this parameter''.

\noindent
\textbf{Common beginner errors to avoid:}
\begin{itemize}\setlength{\itemsep}{2pt}
    \item Passing a positional argument \textit{after} a keyword argument (Python will raise a syntax error).
    \item Providing the same parameter twice (Python will raise ``multiple values for argument ...'').
    \item Forgetting that positional calls must still follow order (keywords remove that requirement).
\end{itemize}

Example call with keyword arguments:
\begin{minted}[fontsize=\small, linenos]{python}
settings, dataset, result = setup_application(
    threshold=2.2,
    print_each=False
)
\end{minted}

\section{Top-Down API Usage}

\subsection{Task 3: Top-Down API Usage from \texttt{main()}}
In this final task, you practice thinking ``top-down'':
\texttt{main()} should read like a clear story of experiments you want to run, while the lower-level functions handle the details.

\noindent
In \texttt{main()}, run at least:
\begin{enumerate}
    \item one default call,
    \item one positional override call,
    \item one keyword override call,
    \item one fit-based call with \texttt{use\_fit=True}.
\end{enumerate}

Print summary metrics for each run so effects are visible.
\noindent
\textbf{Expected outcome:} by changing only the arguments, you can demonstrate different behaviors without changing internal code.
This is exactly what we mean by ``calling your pipeline like an API''.

\subsection{Reference Implementation (Complete Answer)}
The complete Session 5 model answer is provided in:
\begin{itemize}
    \item \texttt{session5/session5\_iris\_model\_answer.py}
\end{itemize}

When reviewing that file, focus on these Session 4-to-5 upgrades:
\begin{itemize}
    \item package import from \texttt{session5.iris\_utils},
    \item optional/default parameters in key functions,
    \item fit-vs-manual threshold selection logic,
    \item multiple call styles demonstrated in \texttt{main()}.
\end{itemize}
\noindent
\textbf{How to use the answer effectively:} do not copy it blindly.
Instead, compare it to your code function-by-function, and explain (in your own words) what each new parameter enables.

\section{Take Home Message}
Session 5 turns your session 4 pipeline into a configurable interface.
\begin{itemize}\setlength{\itemsep}{2pt}
    \item You keep one stable classifier rule.
    \item You add parameters so the same code can run in multiple modes.
    \item You learn a core refactoring habit: change \textit{interfaces} first, keep \textit{logic} stable, and verify behavior through clear runs in \texttt{main()}.
\end{itemize}
